\documentclass[11pt]{article}
\newcommand{\DOCFOOT}{AMPERE POINT — blokada ładowania w wybranych godzinach}
\usepackage{fontspec}
\setmainfont{DejaVu Sans}[Scale=0.86]
\setsansfont{DejaVu Sans}[Scale=0.86]
\usepackage[a4paper,top=2.0cm,bottom=1.8cm,left=1.9cm,right=1.9cm]{geometry}
\usepackage{graphicx}
\usepackage{xcolor}
\usepackage{tikz}
\usetikzlibrary{calc}
\usepackage{enumitem}
\usepackage{fancyhdr}
\usepackage{titlesec}
\usepackage{array}
\usepackage{tabularx}
\usepackage{booktabs}
\usepackage[hidelinks]{hyperref}
\hypersetup{pdfauthor={Dawid Cekała — Ampere Point}, pdfcreator={Ampere Point}}
\usepackage{needspace}

\definecolor{apnavy}{HTML}{1A1A1A}
\definecolor{apnavy2}{HTML}{4A4A4A}
\definecolor{apink}{HTML}{1C2430}
\definecolor{apmuted}{HTML}{5B6675}
\definecolor{apline}{HTML}{E3E3E3}
\definecolor{apbg}{HTML}{F7F7F7}
\definecolor{apgreen}{HTML}{1F8A4C}
\definecolor{apgreenbg}{HTML}{E7F5EC}
\definecolor{apwarn}{HTML}{B9791B}
\definecolor{apwarnbg}{HTML}{FBF1DE}
\definecolor{apaccent}{HTML}{EC6434}

\setlength{\parindent}{0pt}
\setlength{\parskip}{0.35em}
\renewcommand{\arraystretch}{1.25}

\pagestyle{fancy}\fancyhf{}
\renewcommand{\headrulewidth}{0pt}
\renewcommand{\footrulewidth}{0pt}
\fancyfoot[L]{\footnotesize\color{apmuted}\DOCFOOT{} \textperiodcentered{} oprac. Dawid Cekała}
\fancyfoot[R]{\footnotesize\color{apmuted}\thepage}

% --- wordmark (placeholder do czasu otrzymania pliku logo) ---
\newcommand{\apwordmark}[1][1]{\includegraphics[height=\dimexpr #1\dimexpr 0.86cm\relax\relax]{logo_ap.png}}

% --- pasek tytulowy ---
\newcommand{\apheader}[2]{%
\noindent\begin{tikzpicture}
  \node[anchor=north west, inner sep=0pt] at (0,0) {\apwordmark[1.15]};
\end{tikzpicture}\hfill
\raisebox{0.20cm}{\parbox[t]{9.2cm}{\raggedleft\footnotesize\color{apmuted}#2}}\\[0.25cm]
{\color{apaccent}\rule{\textwidth}{1.8pt}}\\[0.30cm]
{\fontspec{DejaVu Sans}\bfseries\color{apnavy}\LARGE #1}\\[0.15cm]}

% --- krok: numer w kolku ---
\newcommand{\kroknum}[1]{%
\tikz[baseline=(k.base)]{\node[circle, fill=apaccent, text=white, inner sep=0pt, minimum size=6.2mm,
  font=\fontspec{DejaVu Sans}\bfseries\small] (k) {#1};}}

% --- ramka informacyjna ---
\newcommand{\apinfo}[2]{%
\par\vspace{0.15cm}
\noindent\begin{tikzpicture}
\node[draw=apline, fill=apbg, line width=0.7pt, rounded corners=5pt, inner sep=8pt,
      text width=\dimexpr\textwidth-2\fboxsep-14pt\relax, align=left]
 {\small\textbf{\color{apnavy}#1}\par\vspace{2pt}\color{apink}#2};
\end{tikzpicture}\par\vspace{0.15cm}}

\newcommand{\apwarnbox}[2]{%
\par\vspace{0.15cm}
\noindent\begin{tikzpicture}
\node[draw=apwarn!45, fill=apwarnbg, line width=0.7pt, rounded corners=5pt, inner sep=8pt,
      text width=\dimexpr\textwidth-2\fboxsep-14pt\relax, align=left]
 {\small\textbf{\color{apwarn}#1}\par\vspace{2pt}\color{apink}#2};
\end{tikzpicture}\par\vspace{0.15cm}}

\newcommand{\apokbox}[2]{%
\par\vspace{0.15cm}
\noindent\begin{tikzpicture}
\node[draw=apgreen!40, fill=apgreenbg, line width=0.7pt, rounded corners=5pt, inner sep=8pt,
      text width=\dimexpr\textwidth-2\fboxsep-14pt\relax, align=left]
 {\small\textbf{\color{apgreen}#1}\par\vspace{2pt}\color{apink}#2};
\end{tikzpicture}\par\vspace{0.15cm}}

% --- komorka kroku: obrazek + opis ---
% #1 numer, #2 tytul, #3 opis, #4 plik
\newcommand{\krok}[4]{%
\begin{minipage}[t]{0.485\textwidth}
  \vspace{0pt}
  \begin{center}\includegraphics[width=\linewidth, height=7.0cm, keepaspectratio]{#4}\end{center}
  \vspace{0.12cm}
  {\kroknum{#1}\ \textbf{\color{apnavy}#2}}\par\vspace{0.05cm}
  {\small\color{apink}#3}
\end{minipage}}

\newcommand{\krokbezobrazka}[3]{%
\begin{minipage}[t]{0.485\textwidth}
  \vspace{0pt}
  \begin{center}
  \begin{tikzpicture}
    \node[draw=apline, fill=apbg, line width=0.7pt, rounded corners=5pt, inner sep=6pt,
          minimum height=7.0cm, text width=5.3cm, align=center]
      {\footnotesize\color{apmuted}(brak zrzutu ekranu\\dla tego kroku)};
  \end{tikzpicture}
  \end{center}
  \vspace{0.12cm}
  {\kroknum{#1}\ \textbf{\color{apnavy}#2}}\par\vspace{0.05cm}
  {\small\color{apink}#3}
\end{minipage}}

% --- os czasu doby: #1 etykieta, #2 lista stref blokady "a/b,c/d"
\newcommand{\osczasu}[2]{%
\begin{tikzpicture}[x=0.62cm, y=1cm]
  \fill[apgreenbg, rounded corners=2pt] (0,0) rectangle (24,0.66);
  \foreach \a/\b in {#2}{ \fill[apwarnbg] (\a,0) rectangle (\b,0.66);
      \draw[apwarn, line width=0.8pt] (\a,0) rectangle (\b,0.66);}
  \draw[apline, line width=0.8pt, rounded corners=2pt] (0,0) rectangle (24,0.66);
  \foreach \h in {0,2,...,24}{ \draw[apline, line width=0.5pt] (\h,0) -- (\h,-0.08);
     \node[below, font=\tiny\color{apmuted}] at (\h,-0.06) {\h};}
  \node[left, font=\small\bfseries\color{apnavy}] at (-0.35,0.33) {#1};
\end{tikzpicture}}

\graphicspath{{screeny/norm/}}

\begin{document}

\apheader{Jak ustawić blokadę ładowania\\[2pt]w wybranych godzinach}{Instrukcja krok po kroku\\ aplikacja Tuya Smart / Smart Life\\ ładowarki AMPERE POINT z Wi-Fi}

\vspace{0.1cm}
{\large\color{apink} Ta instrukcja pokazuje, jak ustawić ładowarkę tak, aby od poniedziałku do piątku nie ładowała samochodu w godzinach 7:00--13:00 oraz 19:00--22:00. Po ustawieniu wystarczy wpiąć wtyczkę o dowolnej porze — resztą zajmie się ładowarka.}

\apokbox{Co osiągniesz}{Samochód można podłączyć o dowolnej porze, bez sprawdzania godzin i bez zmieniania czegokolwiek w aplikacji. Jeżeli podłączenie wypadnie w godzinach blokady, ładowarka sama przerwie ładowanie i wznowi je po zakończeniu tych godzin.}

\vspace{0.2cm}
{\fontspec{DejaVu Sans}\bfseries\color{apnavy}\large Konfiguracja składa się z dwóch części}

\vspace{0.15cm}
\begin{tabularx}{\textwidth}{@{}p{0.5cm}p{4.6cm}X@{}}
\toprule
& \textbf{\color{apnavy}Część} & \textbf{\color{apnavy}Do czego służy}\\
\midrule
A & Harmonogram \par {\footnotesize (4 wpisy godzinowe)} & Wyłącza ładowanie o 7:00 i o 19:00, a włącza je o 13:00 i o 22:00. Działa wtedy, gdy samochód jest już podłączony i ładuje się.\\
B & Automatyzacje \par {\footnotesize (2 reguły)} & Pilnują ładowarki na bieżąco. Jeśli ktoś wepnie wtyczkę w środku godzin blokady, ładowanie zostanie od razu przerwane.\\
\bottomrule
\end{tabularx}

\apinfo{Dlaczego potrzebne są obie części}{Harmonogram wysyła polecenie tylko o wyznaczonej godzinie — nie reaguje na to, co dzieje się pomiędzy. Automatyzacja odwrotnie: reaguje na podłączenie samochodu, ale nie potrafi włączyć ładowania po zakończeniu blokady. Dopiero razem dają pełny efekt.}

\vspace{0.4cm}
{\footnotesize\color{apmuted} Zrzuty ekranu pochodzą z aplikacji w wersji polskiej. Nazwy niektórych pozycji aplikacja wyświetla po angielsku — w instrukcji podajemy je dokładnie tak, jak widać je na ekranie.}

\clearpage

% ============ CZESC A ============
\apheader{Część A — harmonogram}{Cztery wpisy godzinowe\\ poniedziałek--piątek}

Harmonogram ustawiasz w ustawieniach samej ładowarki. Otwórz aplikację, wejdź w swoją ładowarkę, a następnie w Harmonogram (Timer). Dodaj cztery wpisy:

\vspace{0.2cm}
\begin{tabularx}{\textwidth}{@{}p{2.2cm}p{3.2cm}p{3.6cm}X@{}}
\toprule
\textbf{\color{apnavy}Godzina} & \textbf{\color{apnavy}Działanie} & \textbf{\color{apnavy}Powtarzanie} & \textbf{\color{apnavy}Po co}\\
\midrule
07:00 & Wyłącz (OFF) & pon., wt., śr., czw., pt. & Początek droższych godzin rannych\\
13:00 & Włącz (ON) & pon., wt., śr., czw., pt. & Koniec blokady rannej\\
19:00 & Wyłącz (OFF) & pon., wt., śr., czw., pt. & Początek droższych godzin wieczornych\\
22:00 & Włącz (ON) & pon., wt., śr., czw., pt. & Koniec blokady wieczornej\\
\bottomrule
\end{tabularx}

\apwarnbox{Zaznacz powtarzanie w każdym z czterech wpisów}{Wpis bez zaznaczonych dni tygodnia wykona się tylko raz i potem przestanie działać. To najczęstsza przyczyna tego, że harmonogram „działa pierwszego dnia, a potem przestaje”.}

\apinfo{Co z weekendem}{Nie trzeba nic dodawać. W piątek o 22:00 ładowarka zostaje włączona i pozostaje w tym stanie przez sobotę i niedzielę, aż do poniedziałku do godziny 7:00. W weekend ładowanie jest więc możliwe przez całą dobę.}

\vspace{0.4cm}
{\fontspec{DejaVu Sans}\bfseries\color{apnavy}\large Sprawdź jeszcze strefę czasową}

W ustawieniach urządzenia upewnij się, że ustawiona strefa czasowa odpowiada czasowi lokalnemu. Jeśli się nie zgadza, wszystkie godziny w harmonogramie przesuną się o kilka godzin.

\clearpage

% ============ CZESC B ============
\apheader{Część B — automatyzacja pierwsza}{Blokada 7:00--13:00, poniedziałek--piątek\\ kroki 1--4 z 20}

Teraz utworzymy regułę, która pilnuje ładowarki na bieżąco. Jeżeli w godzinach blokady ładowanie z jakiegokolwiek powodu ruszy — na przykład ktoś właśnie wpiął wtyczkę — reguła natychmiast je przerwie.

\vspace{0.35cm}

\krok{1}{Wejdź w zakładkę Smart}{Na dole ekranu wybierz Smart, a następnie u góry zakładkę Automatyzacja. Naciśnij Create (jeśli masz już inne automatyzacje, naciśnij + w prawym górnym rogu).}{s01_automatyzacja_pusta.jpg}
\hfill
\krok{2}{Wybierz rodzaj wyzwalacza}{Z listy wybierz When device status changes — czyli „gdy zmieni się stan urządzenia”.}{s02_wybor_wyzwalacza.jpg}

\vspace{0.42cm}

\krok{3}{Wskaż pojedyncze urządzenie}{Na dole pojawi się okno Device. Naciśnij Select a single device.}{s03_device_single_warunek.jpg}
\hfill
\krok{4}{Wybierz swoją ładowarkę}{Z listy urządzeń wybierz swoją ładowarkę. Na zdjęciu jest to urządzenie o nazwie \emph{timer\_test} — Ty wybierz nazwę własnej ładowarki.}{s04_lista_urzadzen_warunek.jpg}

\clearpage
\apheader{Część B — automatyzacja pierwsza}{ciąg dalszy — kroki 5--8 z 20}

\krok{5}{Wybierz warunek}{Z listy parametrów wybierz Work State, a następnie Charger Charging (stan „ładowarka ładuje”). Zatwierdź przyciskiem Potwierdź w prawym górnym rogu.}{s05_workstate_charging.jpg}
\hfill
\krok{6}{Warunek jest gotowy}{Wracasz do ekranu tworzenia sceny. W sekcji If widnieje teraz wpis Work State : Charger Charging. Teraz trzeba dodać działanie — naciśnij + przy sekcji Then.}{s06_scena_if_gotowy.jpg}

\vspace{0.42cm}

\krok{7}{Wybierz rodzaj działania}{W oknie Dodaj akcję do wykonania wybierz pierwszą pozycję — Device.}{s07_dodaj_akcje.jpg}
\hfill
\krok{8}{Ponownie pojedyncze urządzenie}{Tak samo jak wcześniej: naciśnij Select a single device.}{s08_device_single_akcja.jpg}

\clearpage
\apheader{Część B — automatyzacja pierwsza}{ciąg dalszy — kroki 9--12 z 20}

\krok{9}{Wybierz tę samą ładowarkę}{Z listy wskaż ponownie swoją ładowarkę — tę samą, którą wybrałeś w kroku 4.}{s09_lista_urzadzen_akcja.jpg}
\hfill
\krok{10}{Ustaw działanie: wyłącz}{Wybierz Switch, a pod nim OFF. Zatwierdź przyciskiem Potwierdź.}{s10_switch_off.jpg}

\vspace{0.42cm}

\krok{11}{Sprawdź, co ustawiłeś}{Ekran pokazuje teraz komplet: If — Work State : Charger Charging, Then — Switch : OFF. Zanim zapiszesz, trzeba jeszcze ograniczyć godziny działania reguły. Naciśnij Precondition.}{s11_scena_if_then.jpg}
\hfill
\krok{12}{Otwórz ustawienie godzin}{Naciśnij pozycję Obowiązywanie odcinka czasu — obecnie ustawioną na \emph{Cały dzień}.}{s12_precondition_calydzien.jpg}

\clearpage
\apheader{Część B — automatyzacja pierwsza}{ciąg dalszy — kroki 13--16 z 20}

\krok{13}{Wybierz ustawienie ręczne}{Domyślnie zaznaczone jest Cały dzień. Wybierz ostatnią pozycję — Dostosowane (ręczne ustawienie czasu).}{s13_czas_calydzien.jpg}
\hfill
\krok{14}{Otwórz wybór godzin}{Po zaznaczeniu Dostosowane pojawia się pozycja Okres czasu. Naciśnij ją.}{s14_czas_dostosowane.jpg}

\vspace{0.42cm}

\krok{15}{Ustaw godziny blokady}{Ustaw początek na 06:55 i koniec na 13:00, po czym naciśnij Potwierdź. Pięć minut zapasu przed godziną 7:00 jest celowe — reguła zaczyna czuwać, zanim harmonogram wyłączy ładowanie.}{s15_okres_kola.jpg}
\hfill
\krok{16}{Ustaw dni tygodnia i miasto}{Naciśnij Powtórz i zaznacz od poniedziałku do piątku (odznacz sobotę i niedzielę). Następnie naciśnij Miasto i wskaż swoją miejscowość. Zatwierdź przyciskiem Potwierdź.}{s16_czas_gotowy.jpg}

\clearpage
\apheader{Część B — automatyzacja pierwsza}{zakończenie — kroki 17--20 z 20}

\krok{17}{Wskazanie miasta}{Miasto wybiera się na mapie — wystarczy przybliżyć okolicę swojego adresu i nacisnąć OK. Aplikacja potrzebuje tej informacji, aby prawidłowo liczyć czas lokalny.}{s17_miasto_mapa.jpg}
\hfill
\krok{18}{Zamknij ustawienia godzin}{Przy pozycji Obowiązywanie odcinka czasu widać teraz ustawiony przedział. Naciśnij Zakończ w prawym górnym rogu.}{s18_precondition_gotowy.jpg}

\vspace{0.42cm}

\krok{19}{Nadaj nazwę i zapisz}{Naciśnij Zapisz. Aplikacja poprosi o nazwę reguły — wpisz taką, która od razu mówi, czego dotyczy, na przykład 6:55-13:00 ladowanie\_off. Zatwierdź.}{s19_nazwa_sceny.jpg}
\hfill
\krok{20}{Gotowe}{Reguła pojawia się na liście automatyzacji. Upewnij się, że przełącznik obok niej jest zielony — oznacza to, że reguła jest aktywna.}{s20_lista_gotowa.jpg}

\clearpage

% ============ DRUGA AUTOMATYZACJA ============
\apheader{Część B — automatyzacja druga}{Blokada w godzinach 19:00--22:00\\ poniedziałek--piątek}

Druga reguła jest identyczna jak pierwsza. Różni się wyłącznie godzinami i nazwą.

\vspace{0.2cm}
Wykonaj jeszcze raz wszystkie kroki od 1 do 20, zmieniając tylko dwie rzeczy:

\vspace{0.3cm}
\begin{tabularx}{\textwidth}{@{}p{3.2cm}X@{}}
\toprule
\textbf{\color{apnavy}Krok} & \textbf{\color{apnavy}Co wpisać}\\
\midrule
Krok 15 \par {\footnotesize godziny} & Początek 18:55, koniec 22:00\\
Krok 19 \par {\footnotesize nazwa} & Na przykład 18:55-22:00 ladowanie\_off\\
\bottomrule
\end{tabularx}

\vspace{0.2cm}
Wszystkie pozostałe ustawienia — urządzenie, warunek \emph{Work State : Charger Charging}, działanie \emph{Switch : OFF}, dni od poniedziałku do piątku oraz miasto — pozostają bez zmian.

\apokbox{Po zakończeniu}{Na liście automatyzacji powinny być dwie reguły, obie z zielonym przełącznikiem. W harmonogramie ładowarki powinny być cztery wpisy z zaznaczonymi dniami od poniedziałku do piątku.}

\vspace{0.5cm}
{\fontspec{DejaVu Sans}\bfseries\color{apnavy}\large Jak to działa na co dzień}

\vspace{0.2cm}
\begin{tabularx}{\textwidth}{@{}p{4.4cm}X@{}}
\toprule
\textbf{\color{apnavy}Sytuacja} & \textbf{\color{apnavy}Co się dzieje}\\
\midrule
Wpinasz wtyczkę o 15:00 & Ładowanie rusza normalnie i trwa do 19:00. O 19:00 zostaje przerwane, a o 22:00 wznowione.\\
Wracasz i wpinasz wtyczkę o 19:10 & Ładowanie rusza na chwilę i zostaje samo przerwane przez automatyzację. O 22:00 rusza ponownie, bez Twojego udziału.\\
Wpinasz wtyczkę o 23:00 & Ładowanie rusza normalnie i trwa do rana, do godziny 7:00.\\
Wpinasz wtyczkę w sobotę & Ładowanie działa bez ograniczeń przez cały weekend.\\
\bottomrule
\end{tabularx}

\apinfo{Uwaga o połączeniu z internetem}{Automatyzacje wymagają połączenia ładowarki z internetem. Jeżeli w miejscu ładowania zasięg Wi-Fi bywa słaby, warto rozważyć ustawienie harmonogramu ładowania również w samym samochodzie — wtedy pojazd sam nie pobierze prądu poza wyznaczonymi godzinami, niezależnie od aplikacji.}

\vspace{0.6cm}
{\color{apline}\rule{\textwidth}{0.8pt}}\\[0.15cm]
{\footnotesize\color{apmuted} W razie pytań prosimy o kontakt: hello@amperepoint.com}

\end{document}
